\section{Related Work}
\label{sec:related}

Our work sits at the intersection of active perception, deformable-object manipulation, online system identification, and cross-material simulation benchmarks. The architectural pieces of \PTA{}---a fixed open-loop probing segment, an MLP encoder of the probing trace, and a context-conditioned residual policy---are deliberately minimal and individually well precedented. We position \PTA{} not as a new inference architecture but as a minimal instantiation of active in-episode material identification used to isolate when probing helps or hurts deformable-material manipulation. Below we review each adjacent line of work, emphasizing what \PTA{} \emph{infers}, what it \emph{conditions on}, and what it is \emph{evaluated against}, rather than claiming algorithmic novelty over each.

\paragraph{Active perception for manipulation.}
Robots have long been studied as embodied inference agents that gather information by acting on the world. Early work on active vision~\cite{bajcsy1988active} and exploratory procedures for haptic object recognition~\cite{lederman1987hand} established the principle that specific motions yield specific physical information. In modern learning-based manipulation, tactile exploration has been used to estimate local object geometry and contact properties~\cite{calandra2018more,xu2013tactile}, to inform regrasping policies~\cite{bauza2020tactile}, and to estimate object pose through pushing~\cite{yu2016more,bauza2018data}. These works focus on discrete or local properties of rigid objects---friction coefficients, object poses, stable grasp configurations---and typically treat the estimate as an auxiliary signal rather than as a continuous latent belief that reshapes the downstream motion. \PTA{} extends this line in two ways. First, we target continuous, bulk physical properties of \emph{deformable} media (stiffness, cohesion, plasticity) that cannot be read off from a single contact. Second, we train the belief end-to-end with the downstream task, so the encoder is free to represent whatever aspects of the material response actually matter for the manipulation objective.

\paragraph{Deformable and granular manipulation.}
Manipulating soft and flowable materials has become a focus area as differentiable and particle-based physics engines have matured. Environments such as PlasticineLab~\cite{huang2022plasticinelab} and DiffSkill~\cite{lin2022diffskill} study skill learning on plasticine-like media via differentiable MPM simulation, and RoboCook-style systems~\cite{shi2023robocook,shi2022robocraft} study tool use on dough-like materials. Granular manipulation has been studied both with physics priors~\cite{schenck2017learning} and with learned dynamics models~\cite{matl2020inferring}. A common feature of these works is that training and evaluation use the same material family, or treat material variation via blanket domain randomization. Our benchmark instead organizes evaluation around a \emph{cross-material} split structure---training on the hardest granular material and evaluating zero-shot on cohesive and elastoplastic media---so that the role of material-specific adaptation can be isolated.

\paragraph{Online system identification and adaptation.}
Adapting policies to unseen dynamics has been approached through meta-learning of dynamics and policies~\cite{finn2017model,nagabandi2019learning}, context-conditioned policies that infer a latent task embedding from recent history~\cite{rakelly2019efficient,zintgraf2021varibad}, and rapid online model adaptation~\cite{nagabandi2018deep,clavera2019modelbased}. Robotic sim-to-real pipelines have combined domain randomization with online adaptation of a low-dimensional parameter vector~\cite{yu2017preparing,peng2018sim2real}. These approaches typically assume that the information needed to adapt is present in the course of normal task execution---the policy learns to read its own rollouts. \PTA{} instead \emph{budgets} a dedicated, short probing phase at the start of each episode whose only purpose is to elicit material response. The conditioning latent and the residual policy are otherwise standard. The contribution we attach to this design choice is empirical, not architectural: we use it as a controlled lens to ask when separating ``what material am I holding?'' from ``how do I act?'' helps versus hurts, and our results show that the answer depends on whether probe-induced perturbations relax on the task timescale.

\paragraph{Multi-physics simulation and cross-material benchmarks.}
Recent simulators have made multi-material robotics research practical. Genesis~\cite{genesis2024} provides a unified framework covering rigid body, MPM, SPH, and FEM physics. Isaac Gym~\cite{makoviychuk2021isaac} and Brax~\cite{freeman2021brax} are widely used for GPU-accelerated rigid-body RL but do not natively support deformable media. Benchmarks such as SoftGym~\cite{lin2021softgym} and PlasticineLab~\cite{huang2022plasticinelab} evaluate policies on deformable tasks but typically fix the material family per task. Our benchmark complements these by pairing a single task definition with a deliberately discriminative material sweep (55pp of baseline variation across families), along with parameter-shifted in-family OOD splits and a standardized evaluator. The goal is a narrow, reproducible testbed in which the specific contribution of material-adaptive components---as opposed to task-specific engineering---can be read off the numbers.
